\documentclass[10pt,a4paper]{amsart}
\usepackage[margin=19mm]{geometry}
\usepackage{amsmath,amssymb,mathtools}
\usepackage[scr=rsfs,cal=euler]{mathalfa}
\usepackage{xfrac}
\usepackage[unicode,hidelinks,bookmarks=false]{hyperref}
\newcommand{\ie}{i.e.\ }
\newcommand{\cf}{cf.\ }
\newcommand{\resp}{respectively\ }
\newcommand{\Subsection}{\subsection}
\providecommand{\nobreakdash}{\nobreak\hbox{-}}
\setlength{\emergencystretch}{2em}
\multlinegap0pt
\usepackage{bm}
\usepackage{bbm}
\usepackage{dsfont}
\usepackage{xspace}
\usepackage{cancel}
\usepackage{mathtools}
\usepackage[shortlabels]{enumitem}
\usepackage{amsthm}
\makeatletter
\@ifundefined{theorem}{\newtheorem{theorem}{Theorem}}{}
\@ifundefined{definition}{\newtheorem{definition}[theorem]{Definition}}{}
\@ifundefined{lemma}{\newtheorem{lemma}[theorem]{Lemma}}{}
\@ifundefined{remark}{\newtheorem{remark}[theorem]{Remark}}{}
\makeatother
\makeatletter
\newcommand{\degf}[1]{\left|#1\right|}
\newcommand{\fiz}{\triangleleft}
\newcommand{\fizb}{\blacktriangleleft}
\expandafter\ifx\csname proofw\endcsname\relax
\newenvironment{proofw}{\par
  \pushQED{\qed}%
  \normalfont \topsep6\p@\@plus6\p@\relax
  \trivlist
  \item[]\ignorespaces
}{%
  \popQED\endtrivlist\@endpefalse
}
\fi
\newcommand{\trid}{\triangleright}
\newcommand{\tridb}{\blacktriangleright}
\makeatother
\title[On bicrossed product of fusion categories and exact factoriz]{On bicrossed product of fusion categories and exact factorizations}
\author{Monique M\"uller, H\'ector Mart\'in Pe\~na Pollastri, Julia Plavnik}
\date{Source: arXiv:2405.10207v2 (CC BY 4.0)}
\begin{document}
\maketitle

\noindent\emph{Source and licence.} On bicrossed product of fusion categories and exact factorizations. Version arXiv:2405.10207v2 (CC BY 4.0), distributed under the Creative Commons Attribution 4.0 International licence, as recorded in the arXiv licence metadata for this exact version and shown on the abstract page https://arxiv.org/abs/2405.10207v2. Full rights notice: licenses/arxiv-2405.10207-CC-BY-4.0.txt.\par\smallskip

\noindent\textit{Editorial scope.} Counted unit: the Theorem whose source label is \texttt{thm:bicrossed-product-fusion-rings} in the article (source0.tex lines 698--699, source label \texttt{thm:bicrossed-product-fusion-rings}), delivered together with the article's own complete original proof of it (source0.tex lines 700--749, one \texttt{proof} environment of 50 lines). No mathematics is written, repaired or re-derived: statement and proof are the article's own text line for line. Required context retained uncounted: remark:inversas-matched-pair-grupos (lines 584--586); def:mathced-pair-fusion-rings (lines 641--658); def:mathced-pair-fusion-rings-3 (lines 647--657); def:mathced-pair-fusion-rings-4 (lines 652--655); lemma:dual-of-triangles (lines 661--666). Every definition, display and label of those lines is retained unchanged. Omitted from this package and not redistributed: 1--583, 587--640, 656--660, 667--697, 750--2140. Nothing inside a retained excerpt is omitted or altered: every retained body is delivered exactly as the article's lines are written, so no omission receipt of any kind accompanies this package. Citations: the retained excerpts carry no citation command at all, so no bibliography entry of the article is transcribed and no reference is redistributed. Reference adaptation: none was needed and none was made; every reference inside the delivered unit resolves inside it, so no unresolved reference or citation is produced. Printed slips kept literal: no printed word, symbol, hypothesis or number of the counted statement or of its proof was altered. Layout and class: the article's own class is \texttt{amsart} with packages including geometry, inputenc, latexsym, amsxtra, amscd, ifthen, amsmath, color, multicol, mathdots, fancyvrb, url, amsfonts, amssymb; the wrapper is the lane's pinned \texttt{amsart} editorial wrapper at 10pt on A4, which restates the theorem-like environments the retained text needs and adds the article's own macro definitions that the retained text needs (\texttt{\textbackslash degf}, \texttt{\textbackslash fiz}, \texttt{\textbackslash fizb}, \texttt{\textbackslash trid}, \texttt{\textbackslash tridb}), with extra wrapper packages (bm, bbm, dsfont, xspace, cancel, mathtools, enumitem). The delivered numbering is the unit's own. Assets: no retained span contains a figure environment, a graphics-inclusion command, a PDF-page inclusion, a table or a TikZ picture; no image file is copied into this package, the package declares no asset dependency and no image file is redistributed with it. Licence: version arXiv:2405.10207v2 of this article is distributed under the Creative Commons Attribution 4.0 International licence, as recorded in the arXiv licence metadata for this exact version and shown on the abstract page https://arxiv.org/abs/2405.10207v2. A full rights notice is delivered at licenses/arxiv-2405.10207-CC-BY-4.0.txt. Characters: every retained excerpt is pure ASCII, so the delivered engine, XeLaTeX with its default Latin Modern text font, reports no missing character.
\begin{remark}\label{remark:inversas-matched-pair-grupos}
In a bicrossed product of groups $H\bowtie K$, the inverses are given by $(k\tridb h)^{-1} = (k\fizb h)\tridb h^{-1}$, $(k\fizb h)^{-1} = k^{-1}\fizb (k\tridb h)$. 
\end{remark}

\begin{definition}\label{def:mathced-pair-fusion-rings}
    Let $(\mathsf{A},\mathsf{B}(\mathsf{A}))$, $(\mathsf{C},\mathsf{B}(\mathsf{C}))$ be fusion rings with faithful gradings given by
    \begin{align*}
        \mathsf{A} = \sum_{h\in H} \mathsf{A}_h, && \mathsf{B}(\mathsf{A}) = \bigsqcup_{h\in H} \mathsf{B}(\mathsf{A})_h, &&   \mathsf{C} = \sum_{k\in K} \mathsf{C}_k, && \mathsf{B}(\mathsf{C}) = \bigsqcup_{k\in K} \mathsf{B}(\mathsf{C})_k, 
    \end{align*}
    for groups $H$ and $K$ respectively. A \emph{matched pair of fusion rings} between $\mathsf{A}$ and $\mathsf{C}$ consist of a $8-$tuple $(\mathsf A, \mathsf C, H, K, \tridb, \fizb, \trid, \fiz)$, where
    \begin{enumerate}[leftmargin=*,label=\rm{(\roman*)}]
    \item the tuple $(H, K, \tridb, \fizb)$ is a matched pair between the groups $H$ and $K$,
    \item\label{def:mathced-pair-fusion-rings-2} a $\mathbb{Z}$-linear left action $\trid\colon K\times \mathsf{A} \to \mathsf{A}$ and a $\mathbb{Z}$-linear right action $\fiz\colon \mathsf{C}\times H \to \mathsf{C}$ such that $k \trid \mathsf{B}(\mathsf{
A})_h = \mathsf{B}(\mathsf{A})_{k\tridb h}$ and $\mathsf{B}(\mathsf{C})_k \fiz h= \mathsf{B}(\mathsf{C})_{k \fizb h}$,
    \item\label{def:mathced-pair-fusion-rings-3} for each $k\in K$ and $h\in H$,
    \begin{align*}
        & k\trid (aa') = (k\trid a) ((k\fizb\degf{a})\trid a'), & a,a'\in \mathsf{B} (\mathsf{A}),\\
        &(cc')\fiz h = (c\fiz (\degf{c'}\tridb h)) (c'\fiz h), & c,c'\in\mathsf{B} (\mathsf{C}),
    \end{align*}
 \item\label{def:mathced-pair-fusion-rings-4} for each $k\in K$ and $h\in H$, $k\trid 1=1$ and $1\fiz h=1.$
    \end{enumerate}
    \end{definition}

\begin{lemma}\label{lemma:dual-of-triangles}
Let $(\mathsf A, \mathsf C, H, K, \tridb, \fizb, \trid, \fiz)$ be a matched pair of fusion rings. Then the actions $\trid$ and $\fiz$ satisfy
\begin{align*}
    (k\trid a)^* = (k\fizb \degf{a})\trid a^*, && (c\fiz h)^* = c^*\fiz (\degf{c}\tridb h), && a\in \mathsf{B}(\mathsf{A}), c\in\mathsf{B}(\mathsf{C}), h\in H, k\in K.
\end{align*}
\end{lemma}

\begin{theorem}\label{thm:bicrossed-product-fusion-rings}     The pair $(\mathsf{A}\bowtie \mathsf{C}, \mathsf{B}(\mathsf{A}\bowtie \mathsf{C}))$ with $^*$ is a fusion ring. Furthermore, $(\mathsf{A},\mathsf{B}(\mathsf{A}))$ and $(\mathsf{C},\mathsf{B}(\mathsf{C}))$  are fusion subrings of $\mathsf{A}\bowtie \mathsf{C}$ and the bicrossproduct $\mathsf{A}\bowtie \mathsf{C}$ is an exact factorization of $\mathsf{A}$ and $\mathsf{C}$.
\end{theorem}

\begin{proof} It is not difficult to see that the multiplication defined in $\mathsf{A}\bowtie \mathsf{C}$ is associative by Definition \ref{def:mathced-pair-fusion-rings} \ref{def:mathced-pair-fusion-rings-3} and the structure coefficients are natural numbers given by the formula $N^{a\bowtie c}_{a_1\bowtie c_1, a_2\bowtie c_2}=N^a_{a_1, \degf{c_1}\trid a_2}N^c_{c_1\fiz \degf{a_2}, c_2}$. The unitality condition follows from Definition \ref{def:mathced-pair-fusion-rings} \ref{def:mathced-pair-fusion-rings-4}.

  %       We first show that the multiplication defined in $\mathsf{A}\bowtie \mathsf{C}$ is associative. Note that it is enough to check it for the elements in the basis. On one hand,
  %   \begin{align*}
  %        &((a\bowtie c)(a'\bowtie c'))(a''\bowtie c'') = (a (\degf{c}\trid a') \bowtie (c\fiz \degf{a'}) c')(a''\bowtie c'') \\
  %        &= a (\degf{c}\trid a') (\degf{(c\fiz \degf{a'}) c')}\trid a'')\bowtie (((c\fiz \degf{a'}) c')\fiz \degf{a''} )c''\\
  %        &= a (\degf{c}\trid a') ((\degf{c}\fizb \degf{a'} )\degf{c'})\trid a'')\bowtie ((c\fiz\degf{a'})\fiz (\degf{c'}\tridb \degf{a''}))(c'\fiz \degf{a''})c''\\
  %        &= a (\degf{c}\trid a') ((\degf{c}\fizb \degf{a'}) \degf{c'}\trid a'')\bowtie (c\fiz\degf{a'}(\degf{c'}\tridb \degf{a''}))(c'\fiz \degf{a''})c'',
  %   \end{align*} where in the third equality we used Definition \ref{def:mathced-pair-fusion-rings} \ref{def:mathced-pair-fusion-rings-3}. On the other hand,
  % \begin{align*}
  %        &(a\bowtie c)((a'\bowtie c')(a''\bowtie c'')) = (a\bowtie c)(a'(\degf{c'}\trid a'')\bowtie (c'\fiz \degf{a''})c'')\\
  %        &= a(\degf{c}\trid a'(\degf{c'}\trid a''))\bowtie (c\fiz \degf{a'}(\degf{c'}\tridb \degf{a''}))(c'\fiz \degf{a''})c''\\
  %        &= a(\degf{c}\trid a')((\degf{c}\fizb \degf{a'})\trid (\degf{c'}\trid a''))\bowtie (c\fiz \degf{a'}(\degf{c'}\tridb \degf{a''}))(c'\fiz \degf{a''})c''\\
  %        &= a(\degf{c}\trid a')((\degf{c}\fizb \degf{a'})\degf{c'}\trid a'')\bowtie (c\fiz \degf{a'}(\degf{c'}\tridb \degf{a''}))(c'\fiz \degf{a''})c'',
  %   \end{align*} where in the third equality we used Definition \ref{def:mathced-pair-fusion-rings} \ref{def:mathced-pair-fusion-rings-3}. The unity condition follows from Definition \ref{def:mathced-pair-fusion-rings} \ref{def:mathced-pair-fusion-rings-4}.

    % Since
    % \begin{align}\label{eq:definition-producto-bicrossed-product-anillo-fusion}
    % \begin{aligned}
    %     (a_1\bowtie c_1)(a_2\bowtie c_2) &= a_1 (\degf{c_1}\trid a_2) \bowtie (c_1\fiz \degf{a_2}) c_2 \\
    %     &= \sum_{a\in \mathsf{B}(\mathsf{A})}\sum_{c\in \mathsf{B}(\mathsf{C})}N^a_{a_1, \degf{c_1}\trid a_2}N^c_{c_1\fiz \degf{a_2}, c_2} a\bowtie c,
    %     \end{aligned}
    % \end{align}
    % then $N^{a\bowtie c}_{a_1\bowtie c_1, a_2\bowtie c_2}=N^a_{a_1, \degf{c_1}\trid a_2}N^c_{c_1\fiz \degf{a_2}, c_2}\in \N_0$.

% We will show next that $^*$ is an involution. Indeed, by Lemma \ref{lemma:dual-of-triangles} and Remark \ref{remark:inversas-matched-pair-grupos}, we have
% {\footnotesize\begin{align*}
%     &(a\bowtie c)^{**} = (\degf{c}^{-1}\trid a^* \bowtie c^* \fiz \degf{a}^{-1})^* \\
%     &= \degf{c^* \fiz \degf{a}^{-1}}^{-1} \trid (\degf{c}^{-1}\trid a^*)^* \bowtie (c^* \fiz \degf{a}^{-1})^* \fiz \degf{\degf{c}^{-1}\trid a^*}^{-1}\\
%     &= (\degf{c}^{-1} \fizb \degf{a}^{-1})^{-1} \trid ((\degf{c}^{-1} \fizb \degf{a}^{-1})\trid a^{**} ) \bowtie (c^{**} \fiz (\degf{c}^{-1}\tridb \degf{a}^{-1})) \fiz (\degf{c}^{-1}\tridb \degf{a}^{-1})^{-1}\\
%   &= (\degf{c} \fizb (\degf{c}^{-1}\tridb \degf{a}^{-1})) \trid ((\degf{c}^{-1} \fizb \degf{a}^{-1})\trid a ) \bowtie (c \fiz (\degf{c}^{-1}\tridb \degf{a}^{-1})) \fiz ((\degf{c}^{-1}\fizb \degf{a}^{-1})\tridb \degf{a})\\
%   &= ((\degf{c} \fizb (\degf{c}^{-1}\tridb \degf{a}^{-1}))(\degf{c}^{-1} \fizb \degf{a}^{-1}))\trid a  \bowtie c \fiz( (\degf{c}^{-1}\tridb \degf{a}^{-1}) ((\degf{c}^{-1}\fizb \degf{a}^{-1})\tridb \degf{a}))\\
%   &= ((\degf{c} \degf{c}^{-1})\fizb \degf{a}^{-1})\trid a \bowtie c \fiz (\degf{c}^{-1}\tridb (\degf{a}^{-1} \degf{a}))\\
%   &= (e\fizb \degf{a}^{-1})\trid a \bowtie c \fiz (\degf{c}^{-1}\tridb e) = e\trid a \bowtie c \fiz e = a\bowtie c.
% \end{align*}}

It follows from Lemma \ref{lemma:dual-of-triangles} and Remark \ref{remark:inversas-matched-pair-grupos} that the map $^*$ is indeed an involution. 

We need to check that $N^{1\bowtie 1}_{a_1\bowtie c_1, a_2\bowtie c_2} = 1$ if $a_2\bowtie c_2 = (a_1\bowtie c_1)^* =\degf{c_1}^{-1}\trid a_1^* \bowtie c_1^* \fiz \degf{a_1}^{-1}$ and it is equal $0$ otherwise. It follows that
\begin{align*}
N^{1\bowtie 1}_{a_1\bowtie c_1, a_2\bowtie c_2}=N^1_{a_1, \degf{c_1}\trid a_2}N^1_{c_1\fiz \degf{a_2}, c_2 } = \begin{dcases*}
    1 & if $\degf{c_1}\trid a_2 = a_1^*$ and $c_2= (c_1\fiz \degf{a_2})^*$,\\
    0 & otherwise.
\end{dcases*}
\end{align*}
Then $N^{1\bowtie 1}_{a_1\bowtie c_1, a_2\bowtie c_2} = 1$ when $a_2 =  \degf{c_1}^{-1}\trid a_1^* $ and $c_2 = (c_1\fiz \degf{a_2})^*$ and $0$ otherwise. By Lemma \ref{lemma:dual-of-triangles}, 
%we show that  $(c_1\fiz \degf{a_2})^* =c_1^* \fiz \degf{a_1}^{-1}$. 
we have $(c_1\fiz \degf{a_2})^* = c_1^* \fiz (\degf{c_1}\tridb \degf{a_2})=c_1^* \fiz (\degf{c_1}\tridb (\degf{c_1}^{-1}\tridb \degf{a_1}^{-1})=c_1^* \fiz\degf{a_1}^{-1}$.
  To finish the proof notice that the maps $\iota_{\mathsf A}\colon \mathsf A \to \mathsf A\bowtie \mathsf C$ and $\iota_{\mathsf C}\colon \mathsf C \to \mathsf A\bowtie \mathsf C$ given by $a\mapsto a\bowtie 1$ and $c\mapsto 1\bowtie c$ are inclusions of fusion rings, and $\mathsf A\bowtie \mathsf C =\iota_{\mathsf A}(\mathsf A) \bullet \iota_{\mathsf C}(\mathsf C) \simeq \mathsf A \bullet \mathsf C$.
\end{proof}

\end{document}
